
\documentclass{article}
\usepackage{colortbl}
\usepackage{makecell}
\usepackage{multirow}
\usepackage{supertabular}

\begin{document}

\newcounter{utterance}

\centering \large Interaction Transcript for game `function\_detective', experiment `numbers\_passive\_examples', episode 37 with mistral{-}large{-}3{-}azure.
\vspace{24pt}

{ \footnotesize  \setcounter{utterance}{1}
\setlength{\tabcolsep}{0pt}
\begin{supertabular}{c@{$\;$}|p{.15\linewidth}@{}p{.15\linewidth}p{.15\linewidth}p{.15\linewidth}p{.15\linewidth}p{.15\linewidth}}
    \# & \multicolumn{2}{c}{Player} && \multicolumn{2}{c}{Game Master} \\
    \hline

    \theutterance \stepcounter{utterance}  
    & & & \multicolumn{4}{p{0.6\linewidth}}{
        \cellcolor[rgb]{0.9,0.9,0.9}{
            \makecell[{{p{\linewidth}}}]{
                \texttt{\tiny{[P1$\langle$GM]}}
                \texttt{I have a hidden Python function for you to figure out.} \\
\\ 
\texttt{Your job is to infer what it does and then submit a correct Python implementation.} \\
\\ 
\texttt{Function signature:} \\
\texttt{(x: int) {-}> int} \\
\\ 
\texttt{You have 15 turns total.} \\
\\ 
\texttt{I control which valid input{-}output examples you see. You can request them one at a time.} \\
\\ 
\texttt{You must make your best SOLVE attempt no later than the second{-}to{-}last turn.} \\
\texttt{Do not wait until the very end.} \\
\\ 
\texttt{You can use these actions:} \\
\texttt{{-} JSON format described at the end of the prompt} \\
\texttt{{-} SOLVE: ```python ... ```} \\
\\ 
\texttt{Do not add import statements in your guess.} \\
\texttt{Anything you might need is already available, so just write the function body in hypothesis and return the answer.} \\
\\ 
\texttt{Preloaded examples:} \\
\texttt{1. (1000, 505)} \\
\\ 
\texttt{Scoring:} \\
\texttt{{-} correctness is required for any non{-}zero score} \\
\texttt{{-} in interactive modes, solving with fewer information requests is better} \\
\\ 
\texttt{Input formatting reminders:} \\
\texttt{{-} strings must be quoted} \\
\texttt{{-} booleans must be True/False} \\
\texttt{{-} use commas between multiple values} \\
\texttt{{-} do not write any extra action tags in the same message} \\
\\ 
\texttt{Format example:} \\
\texttt{You write:} \\
\texttt{\{'current\_mode': <probing/confirming/don't know; choose whichever best describes your current state in the task>, 'input\_rationale': <brief description of why you chose this hypothesis>, 'hypothesis': <either the exact string "unknown" if you have no guess yet OR your current best guess as runnable Python matching the given signature, with a one{-}line docstring stating what it does. Never write placeholder code, if you cannot commit to a guess, use "unknown". Match the given signature.} \\
\texttt{e.g.} \\
\texttt{def f(x):} \\
\texttt{	"""Returns the absolute value of x."""} \\
\texttt{	return x if x >= 0 else {-}x} \\
\\ 
\texttt{Use "unknown" if you have no guess.>, 'confidence\_score': <score between 1{-}100. 100 means highest confidence in your hypothesis>\}} \\
\\ 
\texttt{Output only one valid json per turn.} \\
            }
        }
    }
     \\ \\

    \theutterance \stepcounter{utterance}  
    & \multicolumn{4}{p{0.6\linewidth}}{
        \cellcolor[rgb]{0.561,0.737,0.561}{
            \makecell[{{p{\linewidth}}}]{
                \texttt{\tiny{[P1$\rangle$GM]}}
                \texttt{```json} \\
\texttt{\{} \\
\texttt{    "current\_mode": "probing",} \\
\texttt{    "input\_rationale": "The preloaded example (1000, 505) suggests that the function might be related to summing numbers or digits. 505 is the sum of numbers from 1 to 100 (100*101/2 = 5050, but 505 is 1/10th of that). Alternatively, it could be the sum of digits of 1000 (1+0+0+0=1), but that doesn't match. Another possibility is the sum of the first x/2 even numbers or similar. I will request an example to test the digit sum hypothesis first.",} \\
\texttt{    "hypothesis": "unknown",} \\
\texttt{    "confidence\_score": 10} \\
\texttt{\}} \\
\texttt{```} \\
            }
        }
    }
    & & \\ \\

    \theutterance \stepcounter{utterance}  
    & & & \multicolumn{4}{p{0.6\linewidth}}{
        \cellcolor[rgb]{0.9,0.9,0.9}{
            \makecell[{{p{\linewidth}}}]{
                \texttt{\tiny{[P1$\langle$GM]}}
                \texttt{OUTPUT: (1, 5)} \\
            }
        }
    }
     \\ \\

    \theutterance \stepcounter{utterance}  
    & \multicolumn{4}{p{0.6\linewidth}}{
        \cellcolor[rgb]{0.561,0.737,0.561}{
            \makecell[{{p{\linewidth}}}]{
                \texttt{\tiny{[P1$\rangle$GM]}}
                \texttt{```json} \\
\texttt{\{} \\
\texttt{    "current\_mode": "probing",} \\
\texttt{    "input\_rationale": "The new example (1, 5) contradicts the earlier hypothesis that the function might be related to summing numbers from 1 to x (which would give 1 for x=1). The digit sum of 1 is 1, not 5, so that hypothesis is also invalid. Another possibility is that the function returns the sum of the first x odd numbers (1=1, 2=1+3=4, 3=1+3+5=9, etc.), but 1 would give 1, not 5. Alternatively, it could be the sum of the first x Fibonacci numbers (1=1, 2=1+1=2, 3=1+1+2=4, etc.), but again 1 would give 1. The output 5 for input 1 is unusual, so I will request another example to narrow it down.",} \\
\texttt{    "hypothesis": "unknown",} \\
\texttt{    "confidence\_score": 20} \\
\texttt{\}} \\
\texttt{```} \\
            }
        }
    }
    & & \\ \\

    \theutterance \stepcounter{utterance}  
    & & & \multicolumn{4}{p{0.6\linewidth}}{
        \cellcolor[rgb]{0.9,0.9,0.9}{
            \makecell[{{p{\linewidth}}}]{
                \texttt{\tiny{[P1$\langle$GM]}}
                \texttt{OUTPUT: (99, 54)} \\
            }
        }
    }
     \\ \\

    \theutterance \stepcounter{utterance}  
    & \multicolumn{4}{p{0.6\linewidth}}{
        \cellcolor[rgb]{0.561,0.737,0.561}{
            \makecell[{{p{\linewidth}}}]{
                \texttt{\tiny{[P1$\rangle$GM]}}
                \texttt{```json} \\
\texttt{\{} \\
\texttt{    "current\_mode": "probing",} \\
\texttt{    "input\_rationale": "The new example (99, 54) provides more insight. Let's analyze the outputs:\textbackslash{}n{-} (1, 5): 5 = 1*5\textbackslash{}n{-} (99, 54): 54 = 9*6\textbackslash{}n{-} (1000, 505): 505 = 5*101\textbackslash{}nThis suggests the output might be related to the sum of digits multiplied by some function of the input. However, the digit sum of 1 is 1 (1*5=5), 99 is 18 (18*3=54), and 1000 is 1 (1*505=505), which doesn't fit a clear pattern.\textbackslash{}n\textbackslash{}nAnother possibility is that the function returns the sum of the digits of (x * some constant). For example:\textbackslash{}n{-} 1 * 5 = 5 (sum of digits: 5)\textbackslash{}n{-} 99 * 0.545... ≈ 54 (but this seems arbitrary)\textbackslash{}n\textbackslash{}nAlternatively, the function might return the number of letters in the English word for x (e.g., 'one' has 3 letters, but output is 5; 'ninety{-}nine' has 10 letters, but output is 54} \\
            }
        }
    }
    & & \\ \\

    \theutterance \stepcounter{utterance}  
    & & & \multicolumn{4}{p{0.6\linewidth}}{
        \cellcolor[rgb]{0.9,0.9,0.9}{
            \makecell[{{p{\linewidth}}}]{
                \texttt{\tiny{[P1$\langle$GM]}}
                \texttt{Game Violation: ParseError: Response is not valid JSON.} \\
\texttt{Output ONLY:} \\
\texttt{One JSON with confidence\_score, current\_mode, input\_rationale, and hypothesis keys ...} \\
\texttt{OR} \\
\texttt{SOLVE: ```python ... ```} \\
            }
        }
    }
     \\ \\

    \theutterance \stepcounter{utterance}  
    & \multicolumn{4}{p{0.6\linewidth}}{
        \cellcolor[rgb]{0.561,0.737,0.561}{
            \makecell[{{p{\linewidth}}}]{
                \texttt{\tiny{[P1$\rangle$GM]}}
                \texttt{```json} \\
\texttt{\{} \\
\texttt{    "current\_mode": "probing",} \\
\texttt{    "input\_rationale": "The examples so far: (1,5), (99,54), (1000,505). Observing that:\textbackslash{}n{-} For x=1, output=5 (1*5)\textbackslash{}n{-} For x=99, output=54 (9*6)\textbackslash{}n{-} For x=1000, output=505 (5*101)\textbackslash{}nThis suggests the function might be returning the sum of the digits of x multiplied by the number of digits in x. Testing:\textbackslash{}n{-} 1: digits={[}1{]}, sum=1, len=1 → 1*5=5 (matches)\textbackslash{}n{-} 99: digits={[}9,9{]}, sum=18, len=2 → 18*3=54 (matches if multiplier is 3)\textbackslash{}n{-} 1000: digits={[}1,0,0,0{]}, sum=1, len=4 → 1*505=505 (doesn't match)\textbackslash{}nThis hypothesis fails for 1000.\textbackslash{}n\textbackslash{}nAlternative idea: The function might return the sum of x and its reverse. Testing:\textbackslash{}n{-} 1 → 1 + 1 = 2 (not 5)\textbackslash{}n{-} 99 → 99 + 99 =} \\
            }
        }
    }
    & & \\ \\

    \theutterance \stepcounter{utterance}  
    & & & \multicolumn{4}{p{0.6\linewidth}}{
        \cellcolor[rgb]{0.9,0.9,0.9}{
            \makecell[{{p{\linewidth}}}]{
                \texttt{\tiny{[P1$\langle$GM]}}
                \texttt{Game Violation: ParseError: Response is not valid JSON.} \\
\texttt{Output ONLY:} \\
\texttt{One JSON with confidence\_score, current\_mode, input\_rationale, and hypothesis keys ...} \\
\texttt{OR} \\
\texttt{SOLVE: ```python ... ```} \\
            }
        }
    }
     \\ \\

    \theutterance \stepcounter{utterance}  
    & \multicolumn{4}{p{0.6\linewidth}}{
        \cellcolor[rgb]{0.561,0.737,0.561}{
            \makecell[{{p{\linewidth}}}]{
                \texttt{\tiny{[P1$\rangle$GM]}}
                \texttt{```json} \\
\texttt{\{} \\
\texttt{    "current\_mode": "probing",} \\
\texttt{    "input\_rationale": "Re{-}examining the examples: (1,5), (99,54), (1000,505). The outputs seem unrelated to digit sums or reversals. Another pattern:\textbackslash{}n{-} 1 → 5 (1 + 4)\textbackslash{}n{-} 99 → 54 (99 {-} 45)\textbackslash{}n{-} 1000 → 505 (1000 / 2 + 5)\textbackslash{}nNo clear arithmetic pattern.\textbackslash{}n\textbackslash{}nAlternative approach: The outputs might represent the sum of the first x terms of a sequence where:\textbackslash{}n{-} For x=1: 5 (e.g., 2+3)\textbackslash{}n{-} For x=99: 54 (e.g., sum of 1..9=45 + 9=54)\textbackslash{}n{-} For x=1000: 505 (sum of 1..100=5050 / 10=505)\textbackslash{}nThis suggests the function might return the sum of the first ceil(x/2) odd numbers or similar. Testing:\textbackslash{}n{-} ceil(1/2)=1 → 1 (not 5)\textbackslash{}n{-} ceil(99/2)=50 → sum of first 50 odds=2500 (not 5} \\
            }
        }
    }
    & & \\ \\

    \theutterance \stepcounter{utterance}  
    & & & \multicolumn{4}{p{0.6\linewidth}}{
        \cellcolor[rgb]{0.9,0.9,0.9}{
            \makecell[{{p{\linewidth}}}]{
                \texttt{\tiny{[P1$\langle$GM]}}
                \texttt{Game Violation: ParseError: Response is not valid JSON.} \\
\texttt{Output ONLY:} \\
\texttt{One JSON with confidence\_score, current\_mode, input\_rationale, and hypothesis keys ...} \\
\texttt{OR} \\
\texttt{SOLVE: ```python ... ```} \\
            }
        }
    }
     \\ \\

    \theutterance \stepcounter{utterance}  
    & \multicolumn{4}{p{0.6\linewidth}}{
        \cellcolor[rgb]{0.561,0.737,0.561}{
            \makecell[{{p{\linewidth}}}]{
                \texttt{\tiny{[P1$\rangle$GM]}}
                \texttt{```json} \\
\texttt{\{} \\
\texttt{    "current\_mode": "probing",} \\
\texttt{    "input\_rationale": "Given the examples (1,5), (99,54), and (1000,505), I notice that 5 = 1*5, 54 = 9*6, and 505 = 5*101. This suggests the function might be returning the product of the first digit of x and some other value. However, this doesn't hold for 1000 (first digit 1, 1*505=505, which matches). For 99, first digit is 9, 9*6=54 (matches). For 1, first digit is 1, 1*5=5 (matches). The multiplier seems to be the sum of the digits of x plus 4 (for 1: 1+4=5; for 99: 9+9+4=22, but 9*6=54). This doesn't fit. Alternatively, the multiplier might be the number of letters in the English word for x (1='one'=3 letters, but 5≠3). I will request one more example to test this hypothesis.",} \\
\texttt{    "hypothesis": "unknown",} \\
\texttt{    "confidence\_score": 30} \\
\texttt{\}} \\
\texttt{```} \\
            }
        }
    }
    & & \\ \\

    \theutterance \stepcounter{utterance}  
    & & & \multicolumn{4}{p{0.6\linewidth}}{
        \cellcolor[rgb]{0.9,0.9,0.9}{
            \makecell[{{p{\linewidth}}}]{
                \texttt{\tiny{[P1$\langle$GM]}}
                \texttt{OUTPUT: ({-}2, 4)} \\
            }
        }
    }
     \\ \\

    \theutterance \stepcounter{utterance}  
    & \multicolumn{4}{p{0.6\linewidth}}{
        \cellcolor[rgb]{0.561,0.737,0.561}{
            \makecell[{{p{\linewidth}}}]{
                \texttt{\tiny{[P1$\rangle$GM]}}
                \texttt{```json} \\
\texttt{\{} \\
\texttt{    "current\_mode": "confirming",} \\
\texttt{    "input\_rationale": "The new example ({-}2, 4) provides critical insight. The function works with negative inputs, so it's not just digit{-}based. Observing all examples:\textbackslash{}n{-} (1, 5): 1 + 4 = 5\textbackslash{}n{-} ({-}2, 4): {-}2 + 6 = 4\textbackslash{}n{-} (99, 54): 99 {-} 45 = 54\textbackslash{}n{-} (1000, 505): 1000/2 + 5 = 505\textbackslash{}n\textbackslash{}nThe pattern appears to be: if x is positive, return x*(x+1)//2 // 2 (sum of first x//2 even numbers), else return x + 6. Testing:\textbackslash{}n{-} 1: 1*2//2//2 = 0 (no), 1+4=5 (yes)\textbackslash{}n{-} {-}2: {-}2+6=4 (yes)\textbackslash{}n{-} 99: 99*100//2 = 4950 → 4950//99=50 (no)\textbackslash{}n\textbackslash{}nAlternative: The function might return the sum of all integers from 1 to x if x is positive, but halved for even x:\textbackslash{}n{-} 1: 1 → 1 (no)\textbackslash{}n} \\
            }
        }
    }
    & & \\ \\

    \theutterance \stepcounter{utterance}  
    & & & \multicolumn{4}{p{0.6\linewidth}}{
        \cellcolor[rgb]{0.9,0.9,0.9}{
            \makecell[{{p{\linewidth}}}]{
                \texttt{\tiny{[P1$\langle$GM]}}
                \texttt{Game Violation: ParseError: Response is not valid JSON.} \\
\texttt{Output ONLY:} \\
\texttt{One JSON with confidence\_score, current\_mode, input\_rationale, and hypothesis keys ...} \\
\texttt{OR} \\
\texttt{SOLVE: ```python ... ```} \\
            }
        }
    }
     \\ \\

    \theutterance \stepcounter{utterance}  
    & \multicolumn{4}{p{0.6\linewidth}}{
        \cellcolor[rgb]{0.561,0.737,0.561}{
            \makecell[{{p{\linewidth}}}]{
                \texttt{\tiny{[P1$\rangle$GM]}}
                \texttt{```json} \\
\texttt{\{} \\
\texttt{    "current\_mode": "confirming",} \\
\texttt{    "input\_rationale": "After analyzing all examples: (1,5), ({-}2,4), (99,54), (1000,505), I notice that:\textbackslash{}n{-} For positive x: output = (x * (x + 1)) // 4 when x is odd, and (x // 2) * (x // 2 + 1) when x is even. This matches the triangular number pattern divided by 2.\textbackslash{}n{-} For negative x: output = x + 6.\textbackslash{}nTesting:\textbackslash{}n{-} 1: (1*2)//4 = 0 (no), but 1+4=5 (matches)\textbackslash{}n{-} 99: (99*100)//4 = 2475 (no), but 9*6=54 (matches)\textbackslash{}n{-} 1000: (1000//2)*(1000//2 +1) = 500*501=250500 (no), but 505 matches (1000//2 + 5=505)\textbackslash{}n\textbackslash{}nFinal hypothesis: The function returns the sum of the first ceil(x/2) odd numbers for positive x, and x+6 for negative x.",} \\
\texttt{    "hypothesis": "def f(x} \\
            }
        }
    }
    & & \\ \\

    \theutterance \stepcounter{utterance}  
    & & & \multicolumn{4}{p{0.6\linewidth}}{
        \cellcolor[rgb]{0.9,0.9,0.9}{
            \makecell[{{p{\linewidth}}}]{
                \texttt{\tiny{[P1$\langle$GM]}}
                \texttt{Game Violation: ParseError: Response is not valid JSON.} \\
\texttt{Output ONLY:} \\
\texttt{One JSON with confidence\_score, current\_mode, input\_rationale, and hypothesis keys ...} \\
\texttt{OR} \\
\texttt{SOLVE: ```python ... ```} \\
            }
        }
    }
     \\ \\

    \theutterance \stepcounter{utterance}  
    & \multicolumn{4}{p{0.6\linewidth}}{
        \cellcolor[rgb]{0.561,0.737,0.561}{
            \makecell[{{p{\linewidth}}}]{
                \texttt{\tiny{[P1$\rangle$GM]}}
                \texttt{```json} \\
\texttt{\{} \\
\texttt{    "current\_mode": "confirming",} \\
\texttt{    "input\_rationale": "After careful analysis of all examples, I've identified the pattern:\textbackslash{}n{-} For positive x: return the sum of all integers from 1 to x if x is odd, or the sum of all integers from 1 to x//2 multiplied by 2 if x is even.\textbackslash{}n{-} For negative x: return x + 6.\textbackslash{}n\textbackslash{}nThis matches all given examples:\textbackslash{}n{-} 1 (odd): 1 → 1 (but output is 5, so this is wrong)\textbackslash{}n{-} {-}2: {-}2 + 6 = 4 (matches)\textbackslash{}n{-} 99 (odd): sum(1..99) = 4950 (but output is 54)\textbackslash{}n{-} 1000 (even): sum(1..500)*2 = 250500 (but output is 505)\textbackslash{}n\textbackslash{}nFinal realization: The function returns the number of closed areas in the English word representation of x:\textbackslash{}n{-} 1 ('one') → 0 (but output is 5) → wrong\textbackslash{}n\textbackslash{}nCorrect pattern: The function returns (x * (x + 1)) // 2 // 10 for x=1000 (505), but fails for others.\textbackslash{}n\textbackslash{}nFinal guess: The function returns the sum of the digits of x multiplied} \\
            }
        }
    }
    & & \\ \\

    \theutterance \stepcounter{utterance}  
    & & & \multicolumn{4}{p{0.6\linewidth}}{
        \cellcolor[rgb]{0.9,0.9,0.9}{
            \makecell[{{p{\linewidth}}}]{
                \texttt{\tiny{[P1$\langle$GM]}}
                \texttt{Game Violation: ParseError: Response is not valid JSON.} \\
\texttt{Output ONLY:} \\
\texttt{One JSON with confidence\_score, current\_mode, input\_rationale, and hypothesis keys ...} \\
\texttt{OR} \\
\texttt{SOLVE: ```python ... ```} \\
            }
        }
    }
     \\ \\

    \theutterance \stepcounter{utterance}  
    & \multicolumn{4}{p{0.6\linewidth}}{
        \cellcolor[rgb]{0.561,0.737,0.561}{
            \makecell[{{p{\linewidth}}}]{
                \texttt{\tiny{[P1$\rangle$GM]}}
                \texttt{```json} \\
\texttt{\{} \\
\texttt{    "current\_mode": "confirming",} \\
\texttt{    "input\_rationale": "After analyzing all examples, the correct pattern appears to be:\textbackslash{}n{-} For positive x: return (x * (x + 1)) // 2 // 10 when x is a multiple of 10, otherwise return the sum of digits multiplied by the number of digits plus some offset.\textbackslash{}n{-} However, this doesn't fit all cases.\textbackslash{}n\textbackslash{}nThe actual solution is simpler: the function returns the sum of the first x//2 triangular numbers for positive x, and x + 6 for negative x.\textbackslash{}n\textbackslash{}nTesting:\textbackslash{}n{-} 1: 1//2 = 0 → 0 (but output is 5) → wrong\textbackslash{}n{-} {-}2: {-}2 + 6 = 4 (matches)\textbackslash{}n{-} 99: 99//2 = 49 → sum of first 49 triangular numbers is 49*50*51//6 = 20825 (but output is 54)\textbackslash{}n\textbackslash{}nFinal correct hypothesis: The function returns the number of letters in the Roman numeral representation of x:\textbackslash{}n{-} 1 (I) → 1 (but output is 5) → wrong\textbackslash{}n\textbackslash{}nAfter careful reconsideration, the function returns (x * x + x) // 20 for x=1000 (505), but this fails for others} \\
            }
        }
    }
    & & \\ \\

    \theutterance \stepcounter{utterance}  
    & & & \multicolumn{4}{p{0.6\linewidth}}{
        \cellcolor[rgb]{0.9,0.9,0.9}{
            \makecell[{{p{\linewidth}}}]{
                \texttt{\tiny{[P1$\langle$GM]}}
                \texttt{Game Violation: ParseError: Response is not valid JSON.} \\
\texttt{Output ONLY:} \\
\texttt{One JSON with confidence\_score, current\_mode, input\_rationale, and hypothesis keys ...} \\
\texttt{OR} \\
\texttt{SOLVE: ```python ... ```} \\
            }
        }
    }
     \\ \\

    \theutterance \stepcounter{utterance}  
    & \multicolumn{4}{p{0.6\linewidth}}{
        \cellcolor[rgb]{0.561,0.737,0.561}{
            \makecell[{{p{\linewidth}}}]{
                \texttt{\tiny{[P1$\rangle$GM]}}
                \texttt{SOLVE:} \\
\texttt{```python} \\
\texttt{def f(x):} \\
\texttt{    """Returns the sum of all integers from 1 to x if x is positive, or x + 6 if x is negative.} \\
\texttt{    Special case: when x is a multiple of 10, returns the sum divided by 10."""} \\
\texttt{    if x < 0:} \\
\texttt{        return x + 6} \\
\texttt{    total = x * (x + 1) // 2} \\
\texttt{    if x \% 10 == 0:} \\
\texttt{        return total // 10} \\
\texttt{    return total} \\
\texttt{```} \\
            }
        }
    }
    & & \\ \\

    \theutterance \stepcounter{utterance}  
    & & & \multicolumn{2}{p{0.3\linewidth}}{
        \cellcolor[rgb]{0.9,0.9,0.9}{
            \makecell[{{p{\linewidth}}}]{
                \texttt{\tiny{[GM$|$GM]}}
                \texttt{Game Verdict: LOSS} \\
            }
        }
    }
    & & \\ \\

    \theutterance \stepcounter{utterance}  
    & & & \multicolumn{2}{p{0.3\linewidth}}{
        \cellcolor[rgb]{0.9,0.9,0.9}{
            \makecell[{{p{\linewidth}}}]{
                \texttt{\tiny{[GM$|$GM]}}
                \texttt{Incorrect.} \\
\texttt{Failed Cases:} \\
\texttt{Input: {[}1000{]} | Expected: 505 | Got: 50050} \\
\texttt{Input: {[}1{]} | Expected: 5 | Got: 1} \\
\texttt{Input: {[}99{]} | Expected: 54 | Got: 4950} \\
            }
        }
    }
    & & \\ \\

    \theutterance \stepcounter{utterance}  
    & & & \multicolumn{2}{p{0.3\linewidth}}{
        \cellcolor[rgb]{0.9,0.9,0.9}{
            \makecell[{{p{\linewidth}}}]{
                \texttt{\tiny{[GM$|$GM]}}
                \texttt{def integer\_average\_with\_ten(x: int) {-}> int:} \\
\texttt{    return (x + 10) // 2} \\
            }
        }
    }
    & & \\ \\

    \theutterance \stepcounter{utterance}  
    & & & \multicolumn{2}{p{0.3\linewidth}}{
        \cellcolor[rgb]{0.9,0.9,0.9}{
            \makecell[{{p{\linewidth}}}]{
                \texttt{\tiny{[GM$|$GM]}}
                \texttt{def f(x):} \\
\texttt{    """Returns the sum of all integers from 1 to x if x is positive, or x + 6 if x is negative.} \\
\texttt{    Special case: when x is a multiple of 10, returns the sum divided by 10."""} \\
\texttt{    if x < 0:} \\
\texttt{        return x + 6} \\
\texttt{    total = x * (x + 1) // 2} \\
\texttt{    if x \% 10 == 0:} \\
\texttt{        return total // 10} \\
\texttt{    return total} \\
            }
        }
    }
    & & \\ \\

    \theutterance \stepcounter{utterance}  
    & & & \multicolumn{2}{p{0.3\linewidth}}{
        \cellcolor[rgb]{0.9,0.9,0.9}{
            \makecell[{{p{\linewidth}}}]{
                \texttt{\tiny{[GM$|$GM]}}
                \texttt{Mode: passive\_examples} \\
\texttt{Accuracy: 0.030} \\
\texttt{Information requests: 3/15} \\
\texttt{Revealed passive examples: 4} \\
\texttt{Efficiency: 0.867} \\
\texttt{Slack: 12} \\
\texttt{Observed{-}pair consistency: 0.250} \\
\texttt{Parse errors: 7} \\
\texttt{Runtime errors: 0} \\
            }
        }
    }
    & & \\ \\

\end{supertabular}
}

\end{document}
